\documentclass[11pt,a4paper]{article}
\usepackage[a4paper,top=22mm,bottom=23mm,left=23mm,right=23mm]{geometry}
\usepackage{fontspec}
\usepackage{amsmath,amssymb,amsthm,mathtools}
\usepackage{unicode-math}
\setmainfont{DejaVu Serif}
\setsansfont{DejaVu Sans}
\setmonofont{DejaVu Sans Mono}[Scale=0.85]
\setmathfont{Latin Modern Math}
\usepackage{booktabs,longtable,array,enumitem,graphicx,xcolor}
\usepackage{microtype}
\usepackage[unicode=true,colorlinks=true,linkcolor=black,citecolor=black,urlcolor=blue!45!black,bookmarksnumbered=true]{hyperref}
\hypersetup{pdfauthor={Kirtus Justus},pdfsubject={Mathematical models, observation and physical identification},pdfkeywords={relational structures,identifiability,quantum models,calibration,statistical inference}}
\newtheorem{theorem}{Theorem}[section]
\newtheorem{proposition}[theorem]{Proposition}
\newtheorem{lemma}[theorem]{Lemma}
\newtheorem{corollary}[theorem]{Corollary}
\theoremstyle{definition}
\newtheorem{definition}[theorem]{Definition}
\newtheorem{example}[theorem]{Example}
\theoremstyle{remark}
\newtheorem{remark}[theorem]{Remark}
\numberwithin{equation}{section}
\setlength{\parindent}{0pt}
\setlength{\parskip}{5pt plus 1pt}
\setlength{\emergencystretch}{2em}
\setcounter{tocdepth}{2}
\setlist{topsep=4pt,itemsep=2pt,parsep=0pt}
\newcommand{\ProgramTitle}{From Prephysical Hypotheses to Physical Models}
\newcommand{\AuthorLocation}{Port Coquitlam, British Columbia, Canada}

\hypersetup{pdftitle={Reference Modular and Jacobi Estimates for a Half-Filled Fermion Interval},pdfauthor={Kirtus Justus}}
\begin{document}
\thispagestyle{empty}\pagestyle{empty}
{\small\sffamily \ProgramTitle\hfill Paper D7\par}
\vspace*{12mm}
{\LARGE\bfseries\raggedright\hyphenpenalty=10000 Reference Modular and Jacobi Estimates for a Half-Filled Fermion Interval\par}
\vspace{9mm}
{\large Kirtus Justus\par}\vspace{1mm}Theoretical Physicist\par\AuthorLocation\par\vspace{2mm}Baseline edition, 12 November 2026\par\vspace{8mm}
\textbf{Scope}\par
Seven notes on the half-filled reference state of a finite fermion interval and its commuting Jacobi matrix: all-degree Jacobi estimates, exact occupation weights and particle–hole pairing, logarithmic weighted modular control, the sharp logarithmic boundary moment, the sine-kernel endpoint, the fixed-index spectral expansion and the occupation logarithm. Each statement is for the specified reference kernel.\par\vspace{5mm}
\textbf{Contents}\par\vspace{1mm}
{\small\begin{longtable}{@{}r>{\raggedright\arraybackslash}p{0.72\textwidth}r@{}}
\toprule Part & Title & Page\\\midrule\endhead
1 & All-Degree Jacobi Estimates and Conditional Spectral Localization & 3\\
2 & Exact Reference Occupation Weights and Particle-Hole Pairing & 10\\
3 & Logarithmic Weighted Modular Control for the Half-Filled Sine Kernel & 13\\
4 & Sharp Logarithmic Jacobi Boundary Moment for a Half-Filled Fermion Interval & 18\\
5 & Half-Filled Sine-Kernel Endpoint and Extreme Modular Discrepancy & 22\\
6 & Fixed-Index Jacobi Spectral Expansion Through Inverse-Square Order & 29\\
7 & Finite-Interval Occupation Logarithm and Reference-Crossing Bounds & 35\\
\bottomrule\end{longtable}}
\vspace{2mm}\textbf{Formal statements}\par\vspace{1mm}
{\footnotesize Identifiers are those of the program theorem registry. Page numbers refer to this file.\par}
{\footnotesize\begin{longtable}{@{}l>{\raggedright\arraybackslash}p{0.66\textwidth}r@{}}
\toprule Identifier & Statement & Page\\\midrule\endhead
D7-01-01 & Theorem (all-degree alternating-walk bound) & 3\\
D7-02-01 & Exact Reference Occupation Weights and Particle-Hole Pairing & 10\\
D7-03-01 & Theorem 1 & 15\\
D7-03-02 & Corollary 2 (transverse modular correction) & 15\\
D7-03-03 & Lemma 3 (uniform fractional majorant moment) & 16\\
D7-03-04 & Corollary 4 (sufficient sublinear spectral envelope) & 16\\
D7-04-01 & Theorem 1 (sharp second moment) & 18\\
D7-05-01 & Proposition 1 & 22\\
D7-05-02 & Theorem 2 & 23\\
D7-05-03 & Theorem 3 & 24\\
D7-05-04 & Corollary 4 & 25\\
D7-05-05 & Theorem 5 & 25\\
D7-05-06 & Theorem 6 (application of the classical discrete-prolate fixed-index asymptotic) & 26\\
D7-05-07 & Proposition 7 & 26\\
D7-06-01 & Fixed-Index Jacobi Spectral Expansion Through Inverse-Square Order & 29\\
D7-07-01 & Theorem 1 & 36\\
D7-07-02 & Theorem 2 & 36\\
D7-07-03 & Corollary 3 & 36\\
D7-07-04 & Corollary 4 (fully explicit error) & 37\\
D7-07-05 & Corollary 5 & 37\\
\bottomrule\end{longtable}}
\vspace{3mm}\textbf{Acknowledgments and contributions}\par
{\small Kirtus Justus directed the research questions, the selection of mathematical models and the publication standards, is the sole named author, and is responsible for the claims presented here. ChatGPT (OpenAI) assisted with derivations, critical checks and manuscript preparation. Claude (Anthropic), Gemini (Google), Grok (SpaceXAI) and DeepSeek (High-Flyer) contributed to the wider research process; attribution of particular arguments is limited to traceable records. Editorial preparation of this edition: Claude (Anthropic). Assistance does not imply provider endorsement, external peer review or institutional endorsement.\par}
\end{document}